\subsection{Simple Modules}

Recall the following definition (II, §1, No.~10, p.~212).

DEFINITION 1. --- Let A be a ring. An A-module M is called simple if it is nonzero and has no submodule distinct from 0 and M.

An A-module M is simple if and only if M is a simple module over its ring of homotheties \(A_{M}\) . Every simple module is indecomposable, of length 1, and therefore primordial (VIII, p.~31, Proposition 4, b)).

Examples. --- 1) The module \(A_{s}\) is a simple A-module if and only if A is a field (I, §9, No.~1, p.~115, Theorem 1). The simple A-modules are then the vector spaces of dimension 1 over the field A.

\begin{enumerate}
\def\labelenumi{\arabic{enumi})}
\setcounter{enumi}{1}
\item
  Let A be a principal ideal domain (VII, §1, No.~1, p.~1, Definition 1) that is not a field. For every irreducible element \(\pi\) of A, the A-module \(A_s / (\pi)\) is simple, and every simple A-module is isomorphic to such a module (VII, §4, No.~8, p.~25, Remark 4). For \(n \geqslant 2\) , the A-module \(A_s / (\pi^n)\) is indecomposable (VII, §4, No.~8, p.~24, Proposition 8) but not simple.
\item
  Let K be a field, V a nonzero right vector space over the field K, and A a subring of the ring \(\operatorname{End}_{\mathrm{K}}(\mathrm{V})\) containing the endomorphisms of V of finite rank (for example, \(\mathrm{A} = \operatorname{End}_{\mathrm{K}}(\mathrm{V})\) ). Let us prove that V is a simple A-module: let W be a nonzero A-submodule of V and x a nonzero element of W; there exists a linear form \(\varphi\) on V such that \(\varphi(x) \neq 0\) (II, §7, No.~5, p.~300, Theorem 6). For every y in V, the mapping \(z \mapsto y\varphi(z)\) , which is linear of rank \(\leqslant 1\) , belongs to A; we therefore have Ax = V, hence a fortiori W = V, which proves that V is a simple A-module.
\end{enumerate}

PROPOSITION 1. --- Let A be a ring.

\begin{enumerate}
\def\labelenumi{\alph{enumi})}
\item
  Let \(\mathfrak{m}\) be a left ideal of A. The A-module \(\mathrm{A}_s / \mathfrak{m}\) is simple if and only if \(\mathfrak{m}\) is a maximal left ideal.
\item
  Let M be a simple A-module, and let x be a nonzero element of M. We have the equality \(M = Ax\) , the annihilator m of x is a maximal left ideal of A, and the mapping \(a \mapsto ax\) defines, by passing to the quotient, an isomorphism from \(A_{s}/m\) to M.
\item
  Let M be a nonzero A-module. If we have M = Ax for every nonzero element x of M, then M is simple.
\end{enumerate}

The submodules of \(A_{s}/m\) are of the form n/m, where n is a left ideal of A containing m (I, §4, No.~6, p.~41, Theorem 4); consequently, the A-module \(A_{s}/m\) is simple if and only if we have \(m \neq A\) and every left ideal n of A containing m satisfies n/m = 0 or \(n/m = A_{s}/m\) , that is, n = m or n = A. This proves a).

Under the assumptions of b), Ax is a nonzero submodule of M, hence is equal to M. Consequently, the mapping \(a \mapsto ax\) defines, by passing to the quotient, an isomorphism from \(A_{s}/m\) to M; the A-module \(A_{s}/m\) is therefore simple, and the left ideal m is maximal by a). This proves b).

Under the assumptions of c), let N be a nonzero submodule of M. If x is a nonzero element of N, then we have \(Ax \subset N\) and Ax = M, and therefore M = N. Hence M is simple.

COROLLARY 1. --- If the ring A is not reduced to 0, then there exist simple A-modules.

Indeed, by Krull's theorem (I, §8, No.~6, p.~104, Theorem 1), there exist maximal left ideals of A.

COROLLARY 2. --- Let A be a local ring (VIII, p.~26, Definition 1) and r its maximal ideal. The A-module \(A_{s}/r\) is simple, and every simple A-module is isomorphic to \(A_{s}/r\) .

Remarks. --- 1) Let A be a commutative ring and m an ideal of A. Then m is the annihilator (II, §1, No.~12, p.~219, Definition 11) of the A-module \(A_{s}/m\) . Hence, if m and \(m'\) are distinct ideals of A, then the A-modules \(A_{s}/m\) and \(A_{s}/m'\) are not isomorphic. There exists a faithful simple A-module (II, §1, No.~12, p.~219) if and only if (0) is a maximal ideal of A, that is, A is a field.

\begin{enumerate}
\def\labelenumi{\arabic{enumi})}
\setcounter{enumi}{1}
\tightlist
\item
  We can give an example of a noncommutative ring A and two distinct maximal left ideals m and \(m'\) of A such that the A-modules \(A_{s}/m\) and \(A_{s}/m'\) are isomorphic (VIII, p.~52, Exercise 3).
\end{enumerate}

